\documentclass[12pt]{article}
\usepackage[utf8]{inputenc}
\usepackage[T1]{fontenc}
\usepackage{hyperref}
\title{Eudeve Language Resources for K-12 Curriculum Development}
\author{Compiled by Claude Sonnet 5.5 (claude-sonnet-5-5) for Andr\'as Kornai}
\date{}
\begin{document}
\maketitle

\section{Introduction}
Eudeve (ISO 639-3 eud; Glottocode eude1234; also Heve, Dohema) is an extinct Southern Uto-Aztecan (Taracahitic) language of northern Sonora, Mexico, reported extinct since the 1930s and closely related to Opata. Any curriculum would be heritage and revitalization work built on historical documentation. No school programme, community organisation or audio was found. The searches found nothing on contemporary descendant communities.

\section{Primary Educational Resources}
None found. No primers or textbooks for teaching Eudeve exist in the search results.

\section{Linguistic and Reference Materials}
\begin{itemize}
\item Wikipedia, Eudeve language: \url{https://en.wikipedia.org/wiki/Eudeve_language}
\item \emph{Grammatical Sketch of the Heve Language} (translated by Buckingham Smith, 1861; anonymous colonial manuscript), Project Gutenberg: \url{https://gutenberg.org/cache/epub/14419/pg14419-images.html}
\item Endangered Languages Project, Eudeve bibliography (lists, among others, Pennington's 1981 \emph{Arte y Vocabulario de la Lengua Dohema, Heve, o Eudeve}, UNAM; Shaul 1991 on morphosyntax; Lionnet 1986): \url{https://www.endangeredlanguages.com/lang/4927/bibliography}
\item Grambank, Eudeve grammatical features: \url{https://grambank.clld.org/languages/eude1234}
\end{itemize}
The search results mention these works but did not give access to the full text of Pennington, Shaul or Lionnet; check a library for them.

\section{Religious and Community Resources}
None found.

\section{Digital and Audio Resources}
The Gutenberg text above is the only digitized primary source found. No audio exists, since the language is extinct.

\section{Curriculum Development Context}
Teaching material would have to be reconstructed from the colonial grammar and the 20th-century studies cited above. Close attention to Opata and to descendant communities in Sonora would be needed; this report found no information on them.
\end{document}